\section{Global representatives on a small polymer}
\label{subsec:rg-small-polymer-gauge}

The construction in this section follows \cite[\S12.7]{QED3Manuscript}.
The scaled fields belong to the fine domain, but this alone
does not give an analytic gain for their vacuum and ordinary
quadratic kernels. On a small polymer they admit a stronger
representation, with a real gauge, two constant fields and
global fields small enough for an analytic disk of radius
$L^{7/4}$.

\begin{proposition}\label{prop:rg-small-polymer-representatives}
Assume \eqref{eq:rg-contraction-parameters}. Let $X$ be small
and let $(\mathscr A,\mathscr C)\in
\widetilde R_{k+1}\times\mathfrak X_{k+1}$.
There are a real function $\lambda_*$, constant bond
functions $\mathscr K_1,\mathscr K_2$, and global bond
functions $\mathscr A^{\mathrm b},\mathscr C^{\mathrm b}$
on the fine torus such that, on every bond with an endpoint
at distance at most $1$ from $X$,
\begin{equation}\label{eq:rg-small-polymer-gauge-identity}
 \mathscr A_L-d\lambda_*=\mathscr A^{\mathrm b}+\mathscr K_1,
 \qquad
 \mathscr C_L=\mathscr C^{\mathrm b}+\mathscr K_2.
\end{equation}
They satisfy
\begin{equation}\label{eq:rg-small-polymer-global-bounds}
 \begin{split}
 \|\mathscr A^{\mathrm b}\|_\infty
          &\le64L^{-15/8+3\varepsilon/2}\theta_1,\\
 \|\partial\mathscr A^{\mathrm b}\|_\infty
          &\le137L^{-15/8+3\varepsilon/2}\theta_2,\\
 \|\delta_\alpha\partial\mathscr A^{\mathrm b}\|_\infty
          &\le1382L^{-15/8+3\varepsilon/2}\theta_3,\\
 \|\mathscr C^{\mathrm b}\|_\infty&<L^{-7/4}d_1,\qquad
 \|\partial\mathscr C^{\mathrm b}\|_\infty<L^{-7/4}d_2,
 \end{split}
\end{equation}
and
\begin{equation}\label{eq:rg-small-polymer-constant-bounds}
 \|\mathscr K_1\|_\infty<b_1,\qquad
 \|\mathscr K_2\|_\infty<L^{-1+\varepsilon}d_1.
\end{equation}
Consequently
\begin{equation}\label{eq:rg-small-polymer-analytic-disk}
 (t\mathscr A^{\mathrm b},t\mathscr C^{\mathrm b})
                       \in R_k\times\mathfrak X_k
                             \qquad(|t|\le L^{7/4}).
\end{equation}
The construction uses a base point of $X$ and a decomposition
and local real representative of $\mathscr A$, and is
independent of the activity. When $\mathscr C=0$, it gives
$\mathscr C^{\mathrm b}=\mathscr K_2=0$.
\end{proposition}

\begin{proof}
Choose a lift of $X$ and a photon site $c_*$ in it.
Set $r=2LM\ge50$ and $K_{\mathrm{loc}}=K(c_*,2r+5)$, with
sup balls understood through their nonwrapping lifts.
Choose the coarse $M$-cube containing $c_*/L$.
If $\mathscr A=A_0+A_1$ is its domain decomposition, its
real representative $A'$ and gauge $\lambda$ on the
enlarged cube give
\[
 a_0=(A_0)_L=a'_0+d\lambda_L,\qquad
 a'_0=(A')_L,\qquad a_1=(A_1)_L
\]
on $K_{\mathrm{loc}}$. To see that these functions are
defined there, the scaled enlarged cube has side
$(2L+1)LM$ and center $z$ with $|c_*-z|<LM/2$.
The inequality $4LM+6\le L^2M$ puts
$K(c_*,2r+6)$ inside this cube. Moreover
$4r+11<\ell$, since $\ell\ge13L^2M$.
Thus the boxes and all derivative cells used below do
not wrap.

We first record the cutoff estimate with its derivative
support. Put
\[
 f(t)=
 \begin{cases}
 1,&|t|\le1,\\
 1-2(|t|-1)^2,&1\le|t|\le3/2,\\
 2(2-|t|)^2,&3/2\le|t|\le2,\\
 0,&|t|\ge2.
 \end{cases}
\]
For a bond $(c,c+\eta_ke_\mu)$ use the product of the three
functions $f((c_\nu+\eta_k\delta_{\mu\nu}/2-c_{*,\nu})/r)$
as its cutoff multiplier. Denote this operation by
$\operatorname{cut}_r$. The function $f$ is continuously
differentiable, $|f'|\le2$, and $f'$ is Lipschitz with
constant $4$. If a bond function $b$ has local norms
$B_0,B_1,B_\alpha$ on $K_{\mathrm{loc}}$, extension by zero
after multiplication gives
\begin{equation}\label{eq:rg-cutoff-estimate}
 \begin{split}
 \|\operatorname{cut}_r b\|_\infty&\le B_0,\\
 \|\partial\operatorname{cut}_r b\|_\infty
                                      &\le B_1+2B_0/r,\\
 \|\delta_\alpha\partial\operatorname{cut}_r b\|_\infty
                             &\le B_\alpha+12B_1/r+12B_0/r^2.
 \end{split}
\end{equation}
Indeed, the multiplier has oscillation at most
$6|c-c'|/r$, first difference at most $2/r$, and
first-difference oscillation at most $12|c-c'|/r^2$.
Apply the product rule to the first difference.
In its H\"older difference the four terms are bounded
by $B_\alpha |c-c'|^\alpha$, twice
$6B_1|c-c'|/r$, and $12B_0|c-c'|/r^2$.
Here $|c-c'|\le1$ implies
$|c-c'|\le|c-c'|^\alpha$.
If any such term is nonzero, all its cells lie in
$K(c_*,2r+5)$: a nonzero multiplier or its difference
requires a site within $2r+1$, the first-difference
cell adds one mesh length, and the pair adds at most
$1$. The same bound covers the two-step cell at
$k=0$. Thus the local estimates suffice for the global
ones. The cutoff equals one on bonds with both ends in
$K(c_*,r)$.

For $b=a'_0$ or $a_1$, write $v=c-c_*$ and define
\[
 V[b]_{\nu\mu}=(\partial_\nu b_\mu)(c_*),\qquad
 \operatorname{Lin}[b]_\mu(c)=\sum_\nu V[b]_{\nu\mu}v_\nu,
 \qquad
 \operatorname{rem}[b]=b-b(c_*)-\operatorname{Lin}[b].
\]
Its first difference is
$\partial_\nu b_\mu(c)-\partial_\nu b_\mu(c_*)$,
and its H\"older first-difference seminorm is at most
$b_3$. Join the two sites by at most $2r+5$ simultaneous
coordinate steps of length at most $1$. Their difference
is therefore at most $(2r+5)b_3\le3rb_3$. A rectilinear
path to $c_*$ has length at most $3(2r+5)$, and the
remainder vanishes at $c_*$. Keeping the sharper derivative
bound along this path and using $r\ge50$ gives
\[
 \|\operatorname{rem}[b]\|_{K_{\mathrm{loc}}}
              \le3(2r+5)^2b_3\le18r^2b_3,\qquad
 \|\partial\operatorname{rem}[b]\|_{K_{\mathrm{loc}}}
              \le3rb_3.
\]
Equation \eqref{eq:rg-cutoff-estimate} consequently bounds
the three global norms of its cutoff by
$18r^2b_3$, $39rb_3$, and $253b_3$, respectively.

Only a real gauge is available. We remove the constant
and symmetric linear part of $a'_0$, leaving its curvature
and the complex linear part of $a_1$. Set
\[
 V^0=V[a'_0],\qquad
 V^{\mathrm s}=(V^0+(V^0)^T)/2,\qquad
 F=V^0-(V^0)^T,
\]
and use the real quadratic function
\begin{equation}\label{eq:rg-real-quadratic-gauge}
 q_*(c)=\sum_\mu\left[a'_{0,\mu}(c_*)-
                    \frac{\eta_k}{2}V^{\mathrm s}_{\mu\mu}\right]v_\mu
              +\frac12\sum_{\nu,\mu}V^{\mathrm s}_{\nu\mu}v_\nu v_\mu.
\end{equation}
Its finite difference is exactly
$(dq_*)_\mu(c)=a'_{0,\mu}(c_*)+
\sum_\nu V^{\mathrm s}_{\nu\mu}v_\nu$.
The mesh term in \eqref{eq:rg-real-quadratic-gauge} cancels
the diagonal term from differencing the quadratic polynomial.
With $\mathscr K_1=a_1(c_*)$ and
\[
 P_\mu(c)=\sum_\nu
                    \left[\frac12F_{\nu\mu}+V[a_1]_{\nu\mu}\right]v_\nu,
\]
we have the local identity
\begin{equation}\label{eq:rg-local-taylor-gauge}
 \mathscr A_L-d(\lambda_L+q_*)=
       \mathscr K_1+P+\operatorname{rem}[a'_0]+\operatorname{rem}[a_1].
\end{equation}

To cut off $P$ we need a bound that also uses the side
$\ell$. A global mean-zero periodic function whose
oscillation between sites at distance at most $1$ is at
most $B$ has sup norm at most $\ell B$. Any two sites
can be joined by at most $\lceil\ell/2\rceil\le\ell$
simultaneous coordinate steps of length at most $1$;
comparison with the mean gives the claim.
Every $\partial_\nu a_{1,\mu}$ has mean zero, sup norm
at most $b_2$, and this oscillation at most $b_3$.
For $a_0$, the global curl agrees on each derivative
cell with the curl of a local real representative.
Its sup norm is at most $2b_2$, its oscillation at most
$2b_3$, and it has mean zero. The cells of a pair at
distance at most $1$ lie in one enlarged cube, so that
the same representative can be used for their comparison.
It follows that every coefficient of $P$ is bounded by
\begin{equation}\label{eq:rg-linear-slope-bound}
 s=2\min(b_2,\ell b_3).
\end{equation}

Use a second radius
\[
 r'=\min(e_k^{-\varepsilon},\ell/5).
\]
The charge condition in \eqref{eq:rg-contraction-parameters}
gives $r'\ge r$, and $4r'+11<\ell$. Thus its cutoff is
also defined without wrapping. The three norms of $P$
on its support are at most $6.3r's$, $s$, and $0$.
Equation \eqref{eq:rg-cutoff-estimate}, with $r'$ in
place of $r$, bounds the cutoff norms by
$6.3r's$, $13.6s$, and $87.6s/r'$.
If $r'=e_k^{-\varepsilon}$, use $s\le2b_2$;
in units of $(\theta_1,\theta_2,\theta_3)$ these are
at most $(12.6,27.2,175.2)L^{-15/8+\varepsilon}$.
If $r'=\ell/5<e_k^{-\varepsilon}$, use
$s\le2\ell b_3$. Since $\ell e_k^\varepsilon<5$,
the corresponding three bounds are at most
$(63,136,876)L^{-15/8+3\varepsilon/2}$.

Define the global field and real gauge by
\begin{equation}\label{eq:rg-global-representative-construction}
 \begin{split}
 \mathscr A^{\mathrm b}
   &=\operatorname{cut}_r
             \big(\operatorname{rem}[a'_0]+\operatorname{rem}[a_1]\big)
                  +\operatorname{cut}_{r'}P,\\
 \lambda_*&=\lambda_L+q_*\quad\text{on }K_{\mathrm{loc}},
                  \qquad \lambda_*=0\quad\text{elsewhere}.
 \end{split}
\end{equation}
The two remainder cutoffs contribute
$(36r^2,78r,506)b_3$. As $r=2LM$, the charge condition
$160L^{1+3\varepsilon/2}Me_k^\varepsilon\le1$ bounds
the first two by $L^{-15/8+3\varepsilon/2}\theta_1$
and $L^{-15/8+3\varepsilon/2}\theta_2$, respectively.
Adding the linear contribution gives the constants
$64$, $137$, and $1382$ in
\eqref{eq:rg-small-polymer-global-bounds}.

In the external argument take
\[
 \mathscr K_2=\mathscr C_L(c_*),\qquad
 \mathscr C^{\mathrm b}
                   =\operatorname{cut}_r(\mathscr C_L-\mathscr K_2).
\]
A rectilinear path to $c_*$ bounds the undifferentiated
field on $K_{\mathrm{loc}}$ by
$6.3r\|\partial\mathscr C_L\|_\infty$.
The cutoff bounds therefore give
\[
 \|\mathscr C^{\mathrm b}\|_\infty
       \le6.3r\|\partial\mathscr C_L\|_\infty,\qquad
 \|\partial\mathscr C^{\mathrm b}\|_\infty
       \le13.6\|\partial\mathscr C_L\|_\infty.
\]
Use \eqref{eq:rg-scaled-external-thresholds} and
$d_2=e_k^\varepsilon d_1$. The first charge condition
involving $160$ gives the desired bound for the sup norm;
the large-$L$ condition gives
$13.6L^{-2+3\varepsilon/2}\le L^{-7/4}$ for the derivative.
This proves the external bounds and the constant bounds.

The diameter estimate in Appendix~\ref{app:polymer-geometry}
puts $X$ in $K(c_*,M(L+2))$. Every bond with an endpoint
within distance $1$ of $X$ has both ends in $K(c_*,r)$,
since $M(L+2)+2\le2LM$. All cutoffs in the local identity
are one there, and the gauge is $\lambda_L+q_*$ there.
This proves \eqref{eq:rg-small-polymer-gauge-identity}.
Finally, the large-$L$ condition implies
$1382L^{-15/8+3\varepsilon/2}<L^{-7/4}$.
Multiplication by any $|t|\le L^{7/4}$ consequently leaves
all five norms strictly below their domain thresholds,
proving \eqref{eq:rg-small-polymer-analytic-disk}.
\end{proof}

\begin{proposition}\label{prop:rg-small-polymer-kernel-comparison}
Let $F(X)$ be a level-$k$ Grassmann activity holomorphic on
$\widetilde R_k\times\mathfrak X_k$, with the kernel support
and coefficient locality specified in
\S\ref{subsec:rg-activity-spaces}, and covariant under the
independent real gauges of its two connection arguments.
For the fields of Proposition~\ref{prop:rg-small-polymer-representatives},
its kernels satisfy
\begin{equation}\label{eq:rg-small-polymer-kernel-comparison}
 |F_{nm}(X;\mathscr A_L,\mathscr C_L;\xi;\zeta)|
   \le\exp(2(n+m))
        |F_{nm}(X;\mathscr A^{\mathrm b},\mathscr C^{\mathrm b};\xi;\zeta)|.
\end{equation}
For the vacuum kernel the two values are equal.
\end{proposition}

\begin{proof}
For each kernel argument $a$ let $c(a)$ be the photon site
of its ordinary leg, or the first photon site of its pair
leg, and let $\omega(a)\in\{0,1\}$ distinguish unbarred and
barred legs. On kernels the gauge homomorphism has the
form
\begin{equation}\label{eq:rg-kernel-gauge-phase}
 (G_\lambda F)_{nm}(\xi;\zeta)
       =F_{nm}(\xi;\zeta)
          \prod_a\exp\{(-1)^{\omega(a)}ie_k\lambda(c(a))\}.
\end{equation}
This follows by applying the homomorphism to the kernel
expansion; the multiplying product is symmetric in each
set of arguments, and hence preserves their antisymmetry.
The gauge relation \eqref{eq:rg-small-polymer-gauge-identity}
first gives
\[
 G_{\lambda_*}F(X;\mathscr A^{\mathrm b}+\mathscr K_1,
                         \mathscr C^{\mathrm b}+\mathscr K_2)
                       =F(X;\mathscr A_L,\mathscr C_L).
\]
Both pairs lie in the domains: a constant changes only
the sup norm, and
$L^{-7/4}+L^{-7/8+\varepsilon/2}<1$ and
$L^{-7/4}+L^{-1+\varepsilon}<1$.
Every phase from this real gauge has modulus one.

We remove the constant fields by continuation of a real
gauge identity. Write $\mathscr K_j=K_j^{\mathrm R}+iK_j^{\mathrm I}$
with real constant fields. On the connected domain
\[
 \Omega=\{z\in\C^4:|\operatorname{Im}z_1|<2,\,
       |\operatorname{Im}z_2|<2,\,|z_3|<2,\,|z_4|<2\},
\]
set
\[
 A(z)=\mathscr A^{\mathrm b}+z_1K_1^{\mathrm R}+z_2K_1^{\mathrm I},
 \qquad
 C(z)=\mathscr C^{\mathrm b}+z_3K_2^{\mathrm R}+z_4K_2^{\mathrm I},
\]
and set
\[
 \Lambda_z(c)=-
       (z_1K_1^{\mathrm R}+z_2K_1^{\mathrm I}
                  +z_3K_2^{\mathrm R}+z_4K_2^{\mathrm I})\cdot(c-c_*)
                    \quad(c\in K(c_*,r)),
\]
with value zero elsewhere. The first connection belongs to
$\widetilde R_k$: its complex small summand has sup norm
below $(L^{-7/4}+4L^{-7/8+\varepsilon/2})\theta_1<\theta_1$
and derivatives those of $\mathscr A^{\mathrm b}$, and its
remaining constant summand is real. Invariance under real
constant additions handles its unrestricted real part.
The external connection has sup norm below
$(L^{-7/4}+4L^{-1+\varepsilon})d_1<d_1$, and derivative
that of $\mathscr C^{\mathrm b}$. The large-$L$ condition
implies these two strict inequalities.

For fixed kernel arguments the function
\[
 H(z)=F_{nm}(X;A(z),C(z);\xi;\zeta)
             \prod_a\exp\{(-1)^{\omega(a)}ie_k\Lambda_z(c(a))\}
\]
is holomorphic on $\Omega$. When $z$ is real, split
$\Lambda_z$ into its first two and last two terms.
These are independent real gauge functions for the two
connections, with differentials equal to the negatives
of their respective constants on $K(c_*,r)$. Every bond
read by the local gauge condition lies in this ball.
The real gauge identity and
\eqref{eq:rg-kernel-gauge-phase} therefore give
$H(z)=H(0)$ for $z\in\R^2\times(-2,2)^2$.
All derivatives of $H-H(0)$ at zero vanish, so its
holomorphic power series is zero there. The identity
theorem extends this equality to the connected domain
$\Omega$. Evaluating at $z=(1,i,1,i)$ gives the identity
for the two original complex constants.

It remains to bound the inverse phases. By locality,
an ordinary argument lies over $X$. For a pair of the
first kind its first photon site also lies in $X$.
For a boundary pair that site has distance at most
$\eta_k\le1$ from $X$. Hence in all nonzero kernel terms
\[
 |c(a)-c_*|_1\le3(M(L+2)+1)\le9LM.
\]
The constant bounds and the charge conditions give
\[
 \begin{split}
 e_k|K_1^{\mathrm I}\cdot(c(a)-c_*)|
       &\le9LM L^{-7/8+\varepsilon/2}e_k^{1/4+\varepsilon}
                                                    \le1,\\
 e_k|K_2^{\mathrm I}\cdot(c(a)-c_*)|
       &\le9ML^\varepsilon e_k^{2\varepsilon}\le1.
 \end{split}
\]
Each leg phase, and its inverse, thus has modulus at most
$\exp(2)$. Combining this bound with the unit modulus
phases from $\lambda_*$ proves
\eqref{eq:rg-small-polymer-kernel-comparison}.
The vacuum has no leg phases, which proves its equality.
\end{proof}
